\documentclass[11pt]{article}
\usepackage{fontspec}
\setmainfont{Times New Roman}
\setsansfont{Arial}
\setmonofont{Consolas}
\usepackage{amsmath,amssymb}
\usepackage{geometry}
\usepackage{microtype}
\geometry{a4paper,margin=3cm}
\setlength{\parindent}{1.25em}
\setlength{\parskip}{0pt}
\emergencystretch=30em
\allowdisplaybreaks
\raggedbottom
\newcommand{\ideal}[1]{\mathfrak{#1}}
\newcommand{\eqzero}[1]{\equiv 0\,(#1)}

\begin{document}

% Унит: Папер 20. Дирецт Интерславиц Латин транслатион фром те Герман финал-аудитед соурце слице.

\setcounter{footnote}{0}
\section*{20. Алгебраичны критериј абсолутној неразложимости}
\begin{center}
Math. Ann. 85 (1922), стр. 26--33
\end{center}

Полином -- с коефициентами из либовольного, абстрактно определьеного польа -- называје се \emph{абсолутно неразложимы}, ако оставаје неразложимы в алгебраично замкньеном польу, до кторого област коефициентов можно разширити.\footnote{Же всако полье, и то в сушчности једнозначно, можно разширити до алгебраично замкньеного польа, т. ј. до такого, в ктором всакы полином једнеј пременној разпадаје се на линеарне факторы, показал E. Steinitz в својеј \emph{Algebraischen Theorie der Körper} (J. f. M. 137 (1910), стр.~167); тым он дал рационалны еквивалент фундаменталного теорема алгебры.} В следујучем покажем, же абсолутну неразложимост, в противности к неразложимости взходно к заданному польу, можно алгебраично схватити; точнејше важи теорем:

\emph{Всакому полиному од $n\ge2$ пременных с неопредельеными коефициентами можно приредити целу рационалну функцију с целочислеными коефициентами од тых коефициентов и дальных неопредельеных, форму редуцибилности; тако, же за всакы специалны полином једного степена нужно и достаточно условје абсолутној неразложимости јест дано неизчезаньем формы редуцибилности. Другыми словами: за всаку специалну систему вредностиј коефициентов, за ктору форма редуцибилности идентично в неопредельеных изчезаје и ктора не условјује пониженье степена, специалны полином ставаје разложимы, и обратно.}

Ако вместо полиномов сматрјајут се хомогене формы в једнеј пременној више, тада вместо степенового условја выступаје простејше условје, же коефициентна система не изчезаје идентично.

Како непосредно следство теорема добиваје се такоже тврдјенье, најпред поставјено од \mbox{А.~Островски}\footnote{\emph{Zur arithmetischen Theorie der algebraischen Größen}. Gött. Nachr. 1919, стр.~279 (помочна лема, стр.~296). За случај польев составльеных из обычных чисел Островски, колико ми јест знано, такоже дошел до горнього главного теорема; своје изследованьа о том он скоро објави. Же в мојем распореду доказа главны теорем важи за либовольне, абстрактно определьене польа, опираје се на том, же форма редуцибилности образована за неопредельене имаје числове коефициенты независне од специално положено в основу польа.}:

\emph{Всакы абсолутно неразложимы полином с алгебраичными числами како коефициентами оставаје неразложимы модуло всакого простого идеала конечного алгебраичного польа $\Omega$, кторо содржи област коефициентов, изкльучивши највише конечно много простых идеалов из $\Omega$. Особливо $\Omega$ може быти такоже полье рационалных чисел.}

О дальных применьеньах к теорији релативно целых функциј\footnote{D. Hilbert, \emph{Mathematische Probleme}. Gött. Nachr. 1900, стр.~253, проблем 14.} намерјам говорити на другом месту.

1. Најпред покажемо, же \emph{всакы полином од $n\ge2$ пременных с неопредельеными коефициентами јест абсолутно неразложимы}. Именно, положимо:
\begin{equation}
\begin{aligned}
F(x)&=\sum A_{i_1\ldots i_n}x_1^{i_1}\cdots x_n^{i_n} &&(i_1+\cdots+i_n\le l),\\
G(x)&=\sum B_{j_1\ldots j_n}x_1^{j_1}\cdots x_n^{j_n} &&(j_1+\cdots+j_n\le m),\\
H(x)&=F(x)G(x)=\sum C_{k_1\ldots k_n}(A,B)x_1^{k_1}\cdots x_n^{k_n} &&(k_1+\cdots+k_n\le l+m),
\end{aligned}
\tag{1}
\end{equation}
кде $A$ и $B$ значат неопредельене, а $C(A,B)$ ставајут целочислене билинеарне комбинације од $A$ и $B$. Зато квоциенты
\[
D=C(A,B)/C_{0\ldots0}(A,B)
\]
сут рационалне функције квоциентов $A/A_{0\ldots0}$ и $B/B_{0\ldots0}$; але число тых функциј ставаје вечше неж число јих аргументов; бо оно јест одповедајуче
\[
\binom{l+m+n}{n}-1,
\qquad
\binom{l+n}{n}-1,
\qquad
\binom{m+n}{n}-1,
\]
и важи
\[
\binom{l+m+n}{n}+1>
\binom{l+n}{n}+\binom{m+n}{n}
\quad\text{за } n\ge2,\ l>0,\ m>0.
\]
Тако обстаја најмене \emph{једна} алгебраична релација меджу $D(A,B)$ и следовательно такоже меджу $C(A,B)$:
\begin{equation}
 \Phi(Z)\ne0\quad[\text{идентично в }Z_{k_1\ldots k_n}],
 \qquad
 \Phi(C(A,B))=0\quad[\text{идентично в }A,B],
\tag{2}
\end{equation}
кде $\Phi$ значи целочислены полином.

Полином с неопредельеными $Z$ како коефициентами:
\begin{equation}
 E(x)=\sum Z_{k_1\ldots k_n}x_1^{k_1}\cdots x_n^{k_n}
 \qquad(k_1+\cdots+k_n\le l+m)
\tag{3}
\end{equation}
јест зато за $n\ge2$ нужно абсолутно неразложимы.\footnote{Бо неопредельене $Z$ по изразе из 2. не сут нуле од $\mathfrak P$.}

2. Целост тых полиномов $\Phi(Z)$, кторе при субституцији $Z=C(A,B)$ идентично в $A,B$ изчезајут, твори \emph{идеал из полиномов},\footnote{Таке целости обычно называјут се модулы; фактично пак јих карактеризујуче својство -- же оне, поред $\Phi_1$ и $\Phi_2$, содржат такоже $\Phi_1-\Phi_2$, а поред $\Phi$ такоже $\Psi\Phi$, кде $\Psi$ значи либоволны целочислены полином в $Z$ -- јест идеално својство.} точнејше просты идеал $\mathfrak P$; бо ако чрез субституцију изчезаје продукт $\Phi_1\Phi_2$, изчезаје најмене једен фактор. Понеже (2) јест изполньено идентично в $A,B$, оно оставаје в силе за всаку специалну систему вредностиј $A,B$, кде $A,B$ и следовательно такоже $\Gamma=C(A,B)$ значат величины некојего польа. \emph{Коефициенты $\Gamma$ всакого полинома разложимого в разширителном польу зато задовольут условјам $\Phi(\Gamma)=0$, сут нуле од $\mathfrak P$,} или такоже од конечно многих базисных полиномов $\mathfrak P$; ныне иде о доказ обрата.

К тому треба приметити, же все релације меджу $A,B,C$ оставајут, ако полиномы (1) чрез введенье новеј пременној $y$ преобразујут се в хомогене формы $F(x,y)$, $G(x,y)$, $H(x,y)$ степенов $l,m,l+m$. Зато такоже таке коефициенты $\Gamma$ ставајут нулами $\mathfrak P$, при кторых напр. $F(x,y)$\footnote{Полиномы и формы, кторе наставајут, когда неопредельене в $F,G,\ldots,f,g,\ldots$ заменьајут се специалными системами вредностиј, означам повсуду чртоју згора: $\bar F,\bar G,\ldots,\bar f,\bar g,\ldots$.} ставаје степеном од $y$; јим теды чрез преход к нехомогенным полиномам одповедаје пониженье степена од $\bar H(x)$, а не разпаданье. Покаже се, же изкльучивши то пониженье степена обрат всегда важи; особливо, же за \emph{хомогене формы обрат важи без изкльучков, ако не все $\Gamma$ изчезајут; другыми словами, же всака нула од $\mathfrak P$ јест дана параметричным представјеньем $\Gamma=C(A,B)$ с величинами $A,B$ из разширителного польа}.\footnote{Же ``в обче'' нуле од $\mathfrak P$ сут дане параметричным представјеньем, следује из својства простого идеала $\mathfrak P$ определьати неразложиму алгебраичну многовидност. Из того пак, како знано, не следује -- что ја изворно превидела и на что \mbox{А.~Островски} давно обратил моје вниманье -- же параметрично представјенье не може одпасти за жадну специалну систему вредностиј. Тут даны доказ опираје се на елементарнејших појмах.}

То доказанье проведу директным образованьем конечно многих функциј $\Phi(Z)$; именнно чрез Кронецкерову субституцију
\[
 x_i=\xi^{d^{\,n-i}}
\]
приредјам полиномам (1) полиномы од \emph{једнеј} пременној; показујем, же тут наставајут лукы в експонентах, и чрез образованье нормы одповедајучих коефициентов приходим к функцијам $\Phi(Z)$ и к исканному критерију.

3. Нехај буде положено:\footnote{\emph{Festschrift}, стр.~11.}
\begin{equation}
\begin{aligned}
 d&=l+m+1,\\
 i&=i_1d^{n-1}+i_2d^{n-2}+\cdots+i_n,\quad
 j=j_1d^{n-1}+\cdots+j_n,\quad
 k=k_1d^{n-1}+\cdots+k_n.
\end{aligned}
\tag{4}
\end{equation}
Тада чрез Кронецкерову субституцију полиномам $F,G,H$ в (1) одповедајут одповедајуче полиномы
\begin{equation}
\begin{aligned}
 f(\xi)&=\sum A_{i_1\ldots i_n}\xi^i &&(i_1+\cdots+i_n\le l),\\
 g(\xi)&=\sum B_{j_1\ldots j_n}\xi^j &&(j_1+\cdots+j_n\le m),\\
 h(\xi)&=\sum C_{k_1\ldots k_n}(A,B)\xi^k &&(k_1+\cdots+k_n\le l+m).
\end{aligned}
\tag{5}
\end{equation}
То приредјенье јест, како знано, једно-једнозначно, понеже за все в (5) выступајуче индуковане експоненты из $k=0$ следује такоже $k_1=0,\ldots,k_n=0$; и зато из $i+j=i'+j'$ следује такоже $i_1+j_1=i'_1+j'_1;\ldots;i_n+j_n=i'_n+j'_n$. Различне продукты $\xi^i\xi^j$ зато могут дати тен сам експонент $\xi^k$ само тогда, когда и одповедајуче продукты $x_1^{i_1}\cdots x_n^{i_n}x_1^{j_1}\cdots x_n^{j_n}$ ведут к тому самому експоненту. Тако приходит:
\begin{equation}
 h(\xi)=f(\xi)g(\xi),
\tag{6}
\end{equation}
и из изполньености релације (6), тако же в $f(\xi)$ и $g(\xi)$ не выступајут експоненты различне од индукованых, назад следује
\begin{equation}
 H(x)=F(x)G(x).
\tag{7}
\end{equation}
При том \emph{в индукованых експонентах в $f(\xi)$ и $g(\xi)$ в самој ствари выступајут лукы}. Бо највышши експонент од $f(\xi)$ ставаје $ld^{n-1}$, другы највышши $(l-1)d^{n-1}+d^{n-2}$; разлика јест $d^{n-2}(d-1)>1$ заради $d=l+m+1\ge3$, $n\ge2$; и то само важи за $g(\xi)$.

Релације (6) и (7), важече за неопредельене $A,B$, оставајут за всаку специалну систему вредностиј $A,B$. Да бы при том степени остали, хомогенизујем (6) и (7) чрез пременне $\eta$ и $y$. \emph{Хомогена форма $H(x,y)$ зато разпадаје се в алгебраичном разширителном польу тогда и само тогда, когда в њем одповедајуча бинарна форма $h(\xi,\eta)$ може се разложити в два факторы тако, же в факторах не выступајут експоненты различне од индукованых.}

4. Да бы се проведло образованье нормы поменьено на концу 2., нехај
\[
 a_1,b_1;\ldots; a_p,b_p;\quad a_{p+1},b_{p+1};\ldots; a_{p+q},b_{p+q}
\]
значат нове неопредельене. Нехај буде положено
\begin{equation}
\begin{aligned}
 \varphi(\xi,\eta)&=(a_1\xi+b_1\eta)\cdots(a_p\xi+b_p\eta)
     =r_p\xi^p+r_{p-1}\xi^{p-1}\eta+\cdots+r_0\eta^p,\\
 \psi(\xi,\eta)&=(a_{p+1}\xi+b_{p+1}\eta)\cdots(a_{p+q}\xi+b_{p+q}\eta)
     =s_q\xi^q+s_{q-1}\xi^{q-1}\eta+\cdots+s_0\eta^q,\\
 \chi(\xi,\eta)&=\varphi(\xi,\eta)\psi(\xi,\eta)
     =t_{p+q}\xi^{p+q}+t_{p+q-1}\xi^{p+q-1}\eta+\cdots+t_0\eta^{p+q},
\end{aligned}
\tag{8}
\end{equation}
кде $r,s,t$ всегда значат хомогенизоване елементарне симетричне функције; т. ј. те симетричне функције рядов $a,b$, линеарных в всаком јединачном реду, из кторых все целе симетричне функције, кторе сут в всаком јединачном реду хомогене и једнакого степена, складујут се цело и с целочислеными коефициентами -- и то само једным способом.

Ныне образујем с дальными неопредельеными $u_{\mu\nu}$ линеарну форму
\begin{equation}
 \zeta(u,a,b)=\sum u_{\mu\nu}r_\mu s_\nu,
\tag{9}
\end{equation}
где сумованье разширјаје се по заданых индексах $\mu$ и $\nu$. Продукт од $\zeta(u)$ со всими конјугатами, кторе наставајут чрез пермутацију реда $a,b$ и сут различне од $\zeta(u)$ и меджу собоју -- норма $N(\zeta(u))$ од $\zeta(u)$ взходно к польу $t(a,b)$, ако $\zeta(u)$ сматра се како функција польа од $r_\mu(a,b)$ и $s_\nu(a,b)$ -- ставаје потом целочислена симетрична функција всех рядов $a,b$, в всаком реду хомогена и једнакого степена, зато по горньем једнозначно определьена цела целочислена функција од $t(a,b)$ и $u_{\mu\nu}$:
\begin{equation}
 N(\zeta(u,a,b))=T(t(a,b),u)\qquad[\text{идентично в }a,b,u].
\tag{10}
\end{equation}
То само размысльенье показује, же такоже
\[
N(z-\zeta(u))=z^\delta+T_{\delta-1}(t,u)z^{\delta-1}+\cdots+T(t,u)
\qquad[\text{идентично в }a,b,u,z]
\]
где все $T$ значат целе целочислене функције од $t$ и $u$. Обе страны тој идентичности изчезајут за $z=\zeta(u)$; и то, заради алгебраичној независности елементарных симетричных функциј $r(a,b)$, $s(a,b)$, идентично в $r,s$, ако $t(a,b)$ по (8) положи се равно целочисленеј билинеарној комбинацији $w(r,s)$ од $r$ и $s$, а $\zeta(u)$ сматра се како функција од $r$ и $s$. Тако приходит
\begin{equation}
\zeta(u)^\delta+T_{\delta-1}(w(r,s),u)\zeta(u)^{\delta-1}+\cdots+T(w(r,s),u)=0
\quad[\text{идентично в }r,s,u].
\tag{11}
\end{equation}\footnote{(11), ако сумованье в (9) разширјено јест по всех индексах, представја Кронецкерово разширенье Гауссового теорема, и то в сушчности в Кронецкеровом распореду доказа (\emph{Zur Theorie der Formen höherer Stufen}, Berliner Ber. 1883, Werke, 2, str.~417).}

Идентичности (10) и (11) оставајут за все специалне системы вредностиј $\alpha,\beta$ од $a,b$; одповедајуче $\varrho,\sigma$ од $r,s$. Ако $\varrho,\sigma$ сут особливо избране тако, же все продукты $\varrho_\mu\sigma_\nu$ изчезајут -- где $\mu,\nu$ значат индексе выступајуче в (9) -- так же $\zeta(u)$ чрез субституцију изчезаје, тада, ако положимо $w(\varrho,\sigma)=\tau$, по (11) приходит:
\begin{equation}
 T(\tau,u)=0\qquad[\text{идентично в }u].
\tag{12}
\end{equation}

Обратно, нехај $\tau_0,\ldots,\tau_{p+q}$ сут таке, не все изчезајуче системы вредностиј, же (12) јест изполньено. Понеже в алгебраичном разширителном польу обстаја разложенье
\begin{equation}
 \bar\chi(\xi,\eta)=\tau_{p+q}\xi^{p+q}+\cdots+\tau_0\eta^{p+q}
=(\alpha_1\xi+\beta_1\eta)\cdots(\alpha_{p+q}\xi+\beta_{p+q}\eta),
\tag{12'}
\end{equation}
кде в жадном фактору $\alpha$ и $\beta$ не изчезајут једночасно,\footnote{Тој (рационалны) ``фундаменталны теорем алгебры'' јест ексистенцијны теорем, на ктором опираје се безизкльучкова валидност обрата изречена на концу 2.} але некторе $\alpha$ или $\beta$ могут изчезати, ставаје $\tau=t(\alpha,\beta)$. Вставјенье тых вредностиј в (10) показује чрез (12), же $\zeta(u,\alpha,\beta)$ или једен конјугатны фактор изчезаје идентично в $u$. Ако зато, по входној нумерацији од $\alpha,\beta$, ныне положимо $r(\alpha,\beta)=\varrho$, $s(\alpha,\beta)=\sigma$, то значи, же $\bar\chi(\xi,\eta)$ допушчаје најмене једно разложенье:
\begin{equation}
\bar\chi(\xi,\eta)=\bar\varphi(\xi,\eta)\bar\psi(\xi,\eta)
=\sum \varrho_\mu\xi^\mu\eta^{p-\mu}\cdot\sum \sigma_\nu\xi^\nu\eta^{q-\nu},
\tag{13}
\end{equation}
тако, же все продукты $\varrho_\mu\sigma_\nu$ выступајуче в $\zeta(u)$ изчезајут, але не все $\varrho$ или $\sigma$.

5. Степени $p,q$ в (8) нехај ныне будут равне $ld^{n-1}$ и $md^{n-1}$, т. ј. степенам од $f(\xi)$ и $g(\xi)$ в (5). Индексы $\mu,\nu$ разделимо на $\mu^{(1)},\nu^{(1)},\mu^{(2)},\nu^{(2)}$; при том $\mu^{(1)}$ проходи все експоненты, кторе цхыбајут в $f(\xi)$, $\nu^{(1)}$ все експоненты од $\xi$ в $\psi(\xi,\eta)$; и одповедајуче $\nu^{(2)}$ все експоненты, кторе цхыбајут в $g(\xi)$, $\mu^{(2)}$ все експоненты од $\xi$ в $\varphi(\xi,\eta)$. Надалье нехај
\[
 e(\xi)=\sum Z_{k_1\ldots k_n}\xi^k\qquad(k_1+\cdots+k_n\le l+m)
\]
значи полином од једнеј пременној одповедајучи полиному (3); и нехај
\[
 S(Z,u)
\]
значи целочислены полином в $Z$ и $u$, кторы наставаје из $T(t,u)$ -- једнозначно определьене правеј строны (10), кде сумованье в (9) разширјаје се по указаных индексах $\mu^{(1)},\nu^{(1)},\mu^{(2)},\nu^{(2)}$ -- тым, же $t$ заменьајут се коефициентами $Z$ и одповедајучими нулами од $e(\xi)$.

Ако ныне $r,s$ заменьајут се коефициентами $A,B$ и одповедајучими нулами од $f(\xi)$ и $g(\xi)$, тогда чрез ту субституцију $\zeta(u)$ изчезаје при указаном сумованьи. Надалье $t=w(r,s)$ преходи в коефициенты $C(A,B)$ и одповедајуче нуле од $h(\xi)$; зато $T(t,u)$ преходи в $S(C(A,B),u)$. По идентичности (11) зато приходит
\begin{equation}
 S(C(A,B),u)=0\qquad[\text{идентично в }A,B,u].
\tag{14}
\end{equation}
Понеже (14) оставаје за всаку специалну систему вредностиј, коефициенты $\Gamma=C(A,B)$ такых специалных $H(x)$ или обчих $H(x,y)$, кторе в разширителном польу разпадајут се в два факторы, задовольут условју
\begin{equation}
 S(\Gamma,u)=0\qquad[\text{идентично в }u].
\tag{15}
\end{equation}

Ако обратно за, не все изчезајуче, коефициенты $\Gamma$ полинома $\bar H(x)$, одповедајуче формы $\bar H(x,y)$, условје (15) јест изполньено, тогда по горньем то јест једнозначно с $T(\tau,u)=0$, где под $\tau$ разумејут се вси коефициенты од $\bar h(\xi,\eta)$. По (13) зато обстаја в разширителном польу од $\Gamma$ разложенье:
\[
 \bar h(\xi,\eta)=\bar f(\xi,\eta)\cdot\bar g(\xi,\eta),
\]
где в $\bar f$ и $\bar g$ заради указаного сумованьа не выступајут жадне од индукованых различне експоненты; и следовательно обстаја релација
\[
 \bar H(x,y)=\bar F(x,y)\cdot\bar G(x,y).
\]
Из того следује, ако за $y=1$ ниједен фактор не ставаје нултего степена, особливо ако $\Gamma$ не условјујут пониженье степена од $H(x)$, дальна релација
\[
 \bar H(x)=\bar F(x)\cdot\bar G(x).
\]

\emph{Условје (15) јест зато нужно и достаточно за то, же $\bar H(x,y)$, и при изкльученью пониженьа степена такоже $\bar H(x)$, разпадаје се в разширителном польу в два факторы.}

Из того далье следује, же $S(Z,u)$ не може изчезати идентично в $Z$ и $u$, бо в 1. показано јест, же за $n\ge2$ полиномы с неопредельеными како коефициентами, а зато такоже одповедајуче хомогене формы в једнеј пременној више, сут абсолутно неразложиме. Тада, под $U$ разумевајучи степенове продукты од $u$, приходит:
\[
 S(Z,u)=\sum \Phi_\lambda(Z)U_\lambda,
\]
кде $\Phi_\lambda(Z)$ по (14) ставајут полиномами из $\mathfrak P$. Тым јест показано, же $\mathfrak P$ не имаје другых нул неж ты дане параметричным представјеньем $\Gamma=C(A,B)$, где $A$ и $B$ принадлежат алгебраичному разширителному польу од $\Gamma$. Обрат изречены на концу 2. јест тым доказаны.

6. \emph{Условје абсолутној неразложимости} можно ныне непосреднье формуловати. Зато треба само степен од $E(x)$ в (3) разложити на все можне начины в два позитивне члены $l+m$, и к всакому разложеньу образовати форму $S(Z,u)$. Продукт по тых формах
\begin{equation}
 R(Z,u)=\prod S(Z,u)
\tag{16}
\end{equation}
твори тада форму редуцибилности од $E(x)$; бо по горньему всакому разпаданьу од $E(x)$ одповедајуче $E(x,y)$ одповедаје изчезанье једного фактора од (16), и обратно. Тым јест доказаны главны теорем поставјены в уводу и скупај показано, како форма редуцибилности може се в самој ствари поставити в конечно многих кроках.

7. Из ексистенције формы редуцибилности (16) непосреднье следује теорем Островски поменьены в уводу. Нехај
\[
 \bar E(x)=\sum \Gamma_{k_1\ldots k_n}x_1^{k_1}\cdots x_n^{k_n}
 \qquad(k_1+\cdots+k_n\le v)
\]
буде абсолутно неразложимы полином с целыми алгебраичными числами $\Gamma$ како коефициентами; нехај $\Omega$ значи конечно алгебраично числово полье, кторо содржи все $\Gamma$, а $\mathfrak p$ просты идеал из $\Omega$.

Тада форма редуцибилности $R(Z,u)$ образована за точны степен од $\bar E(x)$ при субституцији $Z=\Gamma$ оставаје различна од нулы; ставаје
\[
 R(\Gamma,u)=\sum P_\lambda U_\lambda\ne0\qquad[\text{идентично в }u].
\]
Идеал $\mathfrak r=(P_1,P_2,\ldots)$ јест зато деливы само чрез конечно много простых идеалов из $\Omega$. Ако теды $\mathfrak p$ јест различны од тых конечно многих, тогда $R(\Gamma,u)$ модуло $\mathfrak p$ не изчезаје идентично; и понеже класы остатков модуло $\mathfrak p$ знову творе полье, следује, же $\bar E(x)$ модуло $\mathfrak p$ јест неразложимы, точнејше абсолутно неразложимы. То само важи за $\bar E(x,y)$, также не наставаје ани пониженье степена модуло $\mathfrak p$.

Ако коефициенты од $\bar E(x)$ сут алгебраично дробне числа, тада $\mathfrak p$ треба такоже взяти различны од конечно многих простых идеалов, кторе входјат в обчи именователь $\Lambda$; модуло всех осталых простых идеалов можно заменити $E(x)$ чрез $\Lambda\bar E(x)$, также горне предпоставкы сут изполньене. Тым јест теорем доказаны.

\begin{flushleft}
Ерланген, август 1921.
\end{flushleft}
\begin{center}
(Доставльено 28. 8. 1921.)
\end{center}

\clearpage

\end{document}

